\begin{description}
\item[(a)] A circumpolar star never sets below the horizon. At sea level, the limiting declination is
\[
\delta_m=90^\circ-\varphi=90^\circ-46^\circ53'57''=43^\circ6'3''.
\]\hfill(1 point)
Considering atmospheric refraction $\alpha=34'$, $\delta_m=42^\circ32'3''$. Thus Almach is not circumpolar at sea level. \hfill(1 point)
At 1907\,m the horizon drops, with
\[
\cos\theta=\frac{R}{R+h},\qquad \theta=\arccos\frac{R}{R+h}\approx1.4^\circ.
\]\hfill(2 points)
The source then gives $\delta_m=90^\circ-\varphi-\alpha=41^\circ8'$, so Almach is circumpolar from Mount Ceahl\u{a}u. \hfill(1 point)
\textit{The printed formula omits the horizon-dip term $\theta$. Its numerical result $41^\circ8'$ includes both the $34'$ refraction and the approximately $1.4^\circ$ horizon dip.}
% sic: The printed limiting-declination formula omits theta, though its numerical result includes refraction and the horizon dip.
\item[(b)]
\[
\theta=\frac{1.22\lambda}{D}\quad\Rightarrow\quad D=\frac{1.22\lambda}{\theta}.
\]\hfill(2 points)
With $\theta=10''=4.84\times10^{-5}\,\mathrm{rad}$, $D\approx14\,\mathrm{mm}$. \hfill(1 point)
\item[(c)] Pogson's law gives $E_1/E_2=10^{0.4(m_2-m_1)}$, hence $E_1=10.537E_2$. \hfill(2 points)
The total brightness is $E=E_1+E_2=11.537E_2$. \hfill(1 point)
The printed key writes
\[
\frac{E_1}{E}=2.5^{m-m_1},\qquad m=2.17.
\]\hfill(1 point)
\textit{The base 2.5 in the printed equation is inconsistent with Pogson's law stated above.}
% sic: The exponential base 2.5 in this printed equation is inconsistent with Pogson's law; retained as printed.
\item[(d)]
\[
M=m+5-5\log_{10}d.
\]\hfill(1 point)
$d=355\,\mathrm{ly}=108.896\,\mathrm{pc}$. \hfill(1 point)
Thus $M=-3.04$. \hfill(1 point)
\end{description}
